\chapter{Managing Researchers, Traders and Engineers}\label{ch:fm:managing-researchers-traders-and-engineers}

Three researchers' signals went into one combined forecast that made \$30 million in a year. Asked what each had contributed, the desk measured
what the forecast would have lost without each signal in turn: \$2.5 million, \$2.7 million and \$4.6 million, less than a third of the profit
between them. Two of the signals were two routes to the same effect, so each looked nearly worthless once the other was there. The rest of the
\$30 million had to be split some other way, and the way it was split decided which of the three stayed.

\section{Choosing projects: a research portfolio}

A quantitative desk spends much of its budget on work whose value is uncertain: new signals, new markets, new models, better execution. The
head of desk chooses among projects as an investor chooses among assets, with the same two questions: what is each expected to return, and how do
the returns combine.

\begin{definition}[Research portfolio, project scorecard]\label{def:fm:managing-researchers-traders-and-engineers:portfolio}
A desk's \emph{research portfolio}\index{research portfolio} is the set of research and development projects it funds in a period, chosen under a
budget of money and people. A \emph{project scorecard}\index{project scorecard} records each candidate's cost, probability of success, value if
it succeeds (the present value of the P\&L it adds, net of its running costs), its overlap with what the desk already runs, and the evidence
behind each estimate.
\end{definition}

A project's expected net value is $pV(1-\text{overlap})-c$. The chapter's desk has twelve candidates (illustrative, \$ million): a new futures
signal (cost 1.5, probability 0.30, value 12), options-flow features (2.0, 0.20, 15), an execution model (1.0, 0.60, 4), an alternative-data trial
(2.5, 0.10, 20), a second momentum variant that overlaps the book by 70\% (1.0, 0.40, 6), and seven others. With a budget of \$10 million, choosing
greedily by expected net value per dollar of cost funds six projects with an expected net value of \$8.15 million. Over 20\,000 simulated years
the chosen portfolio returns \$8.08 million on average and loses money in 28.5\% of years; a portfolio chosen at random within the same budget
returns \$3.65 million and loses in 44.8\% (\cref{fig:fm:managing-researchers-traders-and-engineers:portfolio}).

\begin{omfigure}
\begin{tikzpicture}
\begin{axis}[width=10.5cm,height=5.2cm,ybar,bar width=6pt,xlabel={\small realised net value of the year's projects (\$ million)},
  ylabel={\small share of years (\%)},xmin=-13,xmax=41,ymin=0,ymax=20,tick label style={font=\footnotesize},grid=major,grid style={gray!20},
  table/col sep=comma,legend cell align=left,
  legend style={font=\scriptsize,at={(0.5,-0.3)},anchor=north,/tikz/every even column/.append style={column sep=8pt}},legend columns=2]
\addplot[fill=omDef!55,draw=omDef] table[x=centre,y=greedy]{figdata/desk/09-managing-researchers-traders-and-engineers/portfolio.csv};
\addplot[fill=gray!40,draw=gray] table[x=centre,y=random]{figdata/desk/09-managing-researchers-traders-and-engineers/portfolio.csv};
\legend{chosen by expected value per dollar,chosen at random}
\end{axis}
\end{tikzpicture}

\omcaption{The distribution of a year's realised net value from twelve candidate projects and a \$10 million budget, chosen greedily by expected
net value per dollar and at random; 20\,000 simulated years, bins of \$4 million. Most research years lose money on some projects; a good portfolio
loses less often. Data: \texttt{fm\_research.portfolio\_samples}.}\label{fig:fm:managing-researchers-traders-and-engineers:portfolio}
\end{omfigure}

\begin{remark}[The scorecard's weak inputs]\label{rem:fm:managing-researchers-traders-and-engineers:inputs}
The probabilities of success are the weakest numbers on a scorecard and the most important. They come from the desk's own history (how many of
its past projects of each kind reached production, which the research log of One Quant Book 7, chapter 1 records) and they are optimistic when
they come from the proposer. Keeping the history is what turns the portfolio from a debate into an estimate.
\end{remark}

\section{Review without theatre}

A research review (One Quant Book 7, chapter 1) is where a project meets its evidence. Reviews become theatre when their outcome does not depend
on what is presented: when the project's sponsor is senior, when the backtest is shown without the number of variants tried, when a kill criterion
was never written. The procedures that prevent it are already in the series: pre-registration of the hypothesis and of the test, a trial count
kept in the research log, stage gates with kill criteria agreed before the evidence arrives.

\begin{method}[A review that decides]\label{met:fm:managing-researchers-traders-and-engineers:review}
\begin{enumerate}
\item Before the work: record the hypothesis, the test, the kill criterion and the stage gate in the research log.
\item At the review: present the result against the pre-registered test, with the number of trials run and the deflated Sharpe ratio
(One Quant Book 4, chapter 12).
\item Decide by the gate: proceed, stop or change the hypothesis, which starts a new trial count.
\item Review the reviewers once a year: compare their decisions with what reached production and what it earned.
\end{enumerate}
\end{method}

The trial count matters more than any other number in the review. A backtest with a Sharpe ratio of 1.2 over five years of daily data has a
probabilistic Sharpe ratio of 0.996: taken alone, it is almost certainly real. If it is the best of 40 variants tried, the expected best of 40
worthless variants (with a spread of Sharpe ratios of 0.5 a year) is 1.09, and the deflated Sharpe ratio falls to 0.59
(\cref{fig:fm:managing-researchers-traders-and-engineers:dsr}). The same result, after 200 variants, is more likely luck than skill.

\begin{omfigure}
\begin{tikzpicture}
\begin{axis}[width=10.5cm,height=5cm,xmode=log,log basis x=10,xlabel={\small variants tried (log scale)},
  ylabel={\small deflated Sharpe ratio},xmin=1,xmax=500,ymin=0,ymax=1.05,tick label style={font=\footnotesize},grid=major,grid style={gray!20},
  table/col sep=comma,xtick={1,10,100},xticklabels={1,10,100}]
\addplot[omDef,thick,mark=*,mark size=1.3pt] table[x=trials,y=dsr]{figdata/desk/09-managing-researchers-traders-and-engineers/dsr.csv};
\draw[omThm,dashed] (axis cs:1,0.5) -- (axis cs:500,0.5);
\end{axis}
\end{tikzpicture}

\omcaption{The probability that a five-year daily backtest with an annual Sharpe ratio of 1.2 reflects a true Sharpe ratio above the best one
expected by luck among the variants tried, against the number of variants (spread of Sharpe ratios 0.5 a year). Data:
\texttt{fm\_research.dsr\_curve}, on One Quant Book 4's \texttt{firm.multitest}.}\label{fig:fm:managing-researchers-traders-and-engineers:dsr}
\end{omfigure}

\section{Who made the money: attributing credit}

\begin{definition}[Leave-one-out contribution, credit attribution]\label{def:fm:managing-researchers-traders-and-engineers:credit}
A contributor's \emph{leave-one-out contribution}\index{leave-one-out contribution} to a combined result is the result with every contributor
less the result without that contributor. \emph{Credit attribution}\index{credit attribution} is a rule that splits a combined result among
the contributors to it, and on which their recognition and pay are based.
\end{definition}

A combined forecast is a game in the sense of the Shapley value (One Quant Book 12, chapter 6): each set $S$ of contributors has a worth $v(S)$,
the P\&L of the best combination of their signals. Leave-one-out credit, $v(N)-v(N\setminus\{i\})$, measures what each adds last; when two
contributors bring the same information, each adds little last and both look worthless. Order-dependent credit gives each contributor what they
added when they arrived, so the first of two overlapping signals gets everything and the second nothing. The Shapley value averages the order
over all orders.

\begin{proposition}[What the Shapley split guarantees]\label{prop:fm:managing-researchers-traders-and-engineers:shapley}
The Shapley value
$\phi_i=\frac{1}{n!}\sum_{\pi}\bigl(v(P_i^\pi\cup\{i\})-v(P_i^\pi)\bigr)$, the average over orders $\pi$ of what $i$ adds to the set $P_i^\pi$ of
contributors before it, satisfies $\sum_i\phi_i=v(N)$; gives equal credit to contributors who add the same to every set; and gives nothing to a
contributor who adds nothing to any set. Leave-one-out credit satisfies the last two but in general not the first: when contributions are
substitutes, $v(N)-v(N\setminus\{i\})$ summed over $i$ is less than $v(N)$.
\end{proposition}

\begin{proof}[Partial proof]
For each order the added values telescope to $v(N)-v(\emptyset)=v(N)$; averaging preserves the sum. Symmetric contributors have identical
distributions of added values over orders; a contributor who adds nothing adds nothing in every order. For leave-one-out, take two identical
contributors: each adds nothing last, so both get zero while $v(N)>0$. That the Shapley value is the only rule with these properties and
additivity across games is Shapley's theorem (1953), cited in \texttt{omsources}.
\end{proof}

\omcode{code/firm/projsel/firm_projsel.py}{84}{108}{Leave-one-out, order-dependent and Shapley credit on a game $v$.}\label{lst:fm:managing-researchers-traders-and-engineers:shapley}

In the chapter's desk (illustrative: signals A and B correlated at 0.8, C independent, ten years of daily data, the combined P\&L scaled to \$30
million), the rules disagree completely (\cref{fig:fm:managing-researchers-traders-and-engineers:credit}):

\begin{center}
\small
\begin{tabular}{@{}lrrrr@{}}
\toprule
contributor & alone & leave-one-out & order A, B, C & Shapley\\
\midrule
A & 22.77 & 2.46 & 22.77 & 12.63\\
B & 22.84 & 2.67 & 2.59 & 12.77\\
C & 4.55 & 4.64 & 4.64 & 4.61\\
sum & -- & 9.77 & 30.00 & 30.00\\
\bottomrule
\end{tabular}
\end{center}

\begin{omfigure}
\begin{tikzpicture}
\begin{axis}[width=10.5cm,height=5.2cm,ybar,bar width=8pt,ymin=0,ymax=25,ylabel={\small credit (\$ million)},
  xtick={0,1,2},xticklabels={A,B,C},tick label style={font=\footnotesize},table/col sep=comma,enlarge x limits=0.25,
  grid=major,grid style={gray!20},legend cell align=left,
  legend style={font=\scriptsize,at={(0.5,-0.25)},anchor=north,/tikz/every even column/.append style={column sep=8pt}},legend columns=4]
\addplot[fill=gray!40,draw=gray] table[x=k,y=alone]{figdata/desk/09-managing-researchers-traders-and-engineers/credit.csv};
\addplot[fill=omThm!50,draw=omThm] table[x=k,y=loo]{figdata/desk/09-managing-researchers-traders-and-engineers/credit.csv};
\addplot[fill=omProp!55,draw=omProp] table[x=k,y=order_abc]{figdata/desk/09-managing-researchers-traders-and-engineers/credit.csv};
\addplot[fill=omDef!55,draw=omDef] table[x=k,y=shapley]{figdata/desk/09-managing-researchers-traders-and-engineers/credit.csv};
\legend{alone,leave-one-out,{order A, B, C},Shapley}
\end{axis}
\end{tikzpicture}

\omcaption{Credit for a combined forecast's \$30 million among three contributors under four rules: each signal alone, leave-one-out,
order of arrival (A first) and the Shapley value. A and B overlap (correlation 0.8), C is independent. Illustrative data. Data:
\texttt{fm\_research.credit}.}\label{fig:fm:managing-researchers-traders-and-engineers:credit}
\end{omfigure}

A bonus pool of \$6 million split in proportion to each rule's credits gives A, B and C \$1.51, \$1.64 and \$2.85 million by leave-one-out, \$2.53,
\$2.55 and \$0.92 million by Shapley, and \$4.55, \$0.52 and \$0.93 million by order of arrival. Leave-one-out pays C most, because C is the only
one whose signal cannot be replaced; order of arrival pays whoever came first. Shapley pays A and B as a pair for the effect they both found, and C
for what only C found. The credits are themselves estimates: rerun on another ten years of data (exercise 7), the split moves by several million.

\section{Keeping research honest}

Credit rules shape research as much as they measure it. Paid by leave-one-out, researchers avoid overlapping work, including the replication
that would expose a spurious signal; paid by order of arrival, they race to be first and hoard; paid by stand-alone P\&L, they rediscover what the
book already holds. Paid by Shapley, a researcher who finds a second route to an effect the desk already exploits is paid for part of it, which
rewards robustness and replication. Three further rules keep research honest whatever the credit rule: every trial goes in the research log;
results are presented against a pre-registered test; and the people who review a signal are not those who are paid for it.

\section{Traders and engineers: different work, different feedback}

Researchers, traders and engineers get feedback at different speeds, and managing them the same way fails all three
(\cref{fig:fm:managing-researchers-traders-and-engineers:loops}). A trader sees the market's verdict in minutes and must be judged on decisions,
not only on the day's P\&L, whose noise dominates it: attribution of P\&L to risk taken and to execution (One Quant Book 1, chapter 7) is the
tool. A researcher's work pays off, if ever, after months of testing and years of live trading: the measure is the quality of the process, trials
and evidence, until the P\&L can speak. An engineer's work shows in systems that run or do not: delivery, reliability and the time it takes to
restore service are measured daily (One Quant Book 15, chapter 28), and the value of the work is the trading it enables, which is why engineers
belong in the \omterm{def:fm:managing-researchers-traders-and-engineers:credit}{credit attribution} of the strategies they make possible.

\begin{omfigure}
\begin{tikzpicture}[font=\footnotesize,>=stealth,box/.style={draw,rounded corners,minimum height=0.62cm,align=center,fill=omDef!10,
  minimum width=3.1cm}]
\node[box] (t) at (0,2.4) {trader\\{\scriptsize decision $\to$ market: minutes}};
\node[box] (e) at (0,1.2) {engineer\\{\scriptsize change $\to$ production: days}};
\node[box] (r) at (0,0) {researcher\\{\scriptsize idea $\to$ live P\&L: months to years}};
\node[align=left,anchor=west,font=\scriptsize] at (2.8,2.4) {judge decisions against attribution;\\one day's P\&L is mostly noise};
\node[align=left,anchor=west,font=\scriptsize] at (2.8,1.2) {judge delivery and reliability;\\share credit for the trading enabled};
\node[align=left,anchor=west,font=\scriptsize] at (2.8,0) {judge process and evidence (trials, tests)\\until the P\&L can speak};
\end{tikzpicture}

\omcaption{Three kinds of work and the speed of their feedback: how each should be judged while its results are still noise.}\label{fig:fm:managing-researchers-traders-and-engineers:loops}
\end{omfigure}

\section{Tutorial: a year of research, and the credit for it}\label{tut:fm:managing-researchers-traders-and-engineers:tut}

\begin{tutorial}
\textbf{Goal.} Choose a \omterm{def:fm:managing-researchers-traders-and-engineers:portfolio}{research portfolio}, deflate a result by its trials, and split the credit for a combined forecast. \textbf{End state:} the
credit table and \cref{fig:fm:managing-researchers-traders-and-engineers:credit}.
\begin{tutsteps}
\item \textbf{The portfolio.} \texttt{firm.projsel.select(fm\_research.PROJECTS, 10)} chooses six projects; \texttt{fm\_research.portfolio()}
simulates 20\,000 years of the choice and of random choices.
\item \textbf{The trials.} \texttt{fm\_research.deflated(trials=40)} gives the probabilistic and deflated Sharpe ratios of a five-year backtest
with a Sharpe ratio of 1.2.
\item \textbf{The game.} \texttt{firm.projsel.game(signals, target)} builds $v$; \texttt{leave\_one\_out}, \texttt{ordered} and \texttt{shapley}
split it (\cref{lst:fm:managing-researchers-traders-and-engineers:shapley}).
\item \textbf{The pool.} \texttt{fm\_research.pool\_split()} turns each rule's credit into shares of a \$6 million pool.
\end{tutsteps}
\end{tutorial}

\textbf{What to change next.} Make A and B identical and check that Shapley pays them equally and leave-one-out pays them nothing; rerun on
another seed and measure how uncertain each credit is.

\section{Build: research portfolio and credit}\label{bld:fm:managing-researchers-traders-and-engineers:projsel}

\begin{build}
\bfield{Purpose} The head of desk's two allocation problems: which projects to fund, and how to credit the people whose work combined.
\bfield{Interface} \texttt{firm.projsel}: \texttt{Project(name, cost, p, value, overlap)}, \texttt{expected\_net}; \texttt{select(projects,
budget)}, \texttt{select\_random}, \texttt{simulate\_year}; \texttt{game(signals, target)}; \texttt{leave\_one\_out}, \texttt{ordered},
\texttt{shapley(v, n, samples, rng)}. The chapter uses One Quant Book 4's \texttt{firm.multitest} for deflated Sharpe ratios.
\bfield{Rules} Projects with non-positive expected net value are never funded; the game's worth is the in-sample P\&L of the least-squares
combination; exact Shapley values for up to about eight contributors, sampled beyond.
\bfield{Acceptance tests} \texttt{code/firm/projsel/tests/}: greedy selection within the budget; Shapley efficiency, symmetry of identical
contributors and zero for a dummy; leave-one-out zero for a duplicated contributor; sampled Shapley close to exact.
\bfield{Stretch} Out-of-sample games (the worth of $S$ measured on data not used to fit the weights); a portfolio with correlated project
outcomes; credit with a time dimension (who kept a signal alive).
\end{build}

\begin{omsources}
\item L.~S.~Shapley, ``A value for n-person games'', in \emph{Contributions to the Theory of Games II}, 1953.
\item S.~Lipovetsky and M.~Conklin, ``Analysis of regression in game theory approach'', \emph{Applied Stochastic Models in Business and Industry}
17(4), 2001.
\item D.~H.~Bailey and M.~López de Prado, ``The deflated Sharpe ratio'', \emph{Journal of Portfolio Management} 40(5), 2014.
\item C.~R.~Harvey, Y.~Liu and H.~Zhu, ``\dots\ and the cross-section of expected returns'', \emph{Review of Financial Studies} 29(1), 2016.
\end{omsources}

\section{Exercises}

\begin{exercise}[$\star$]\label{exo:fm:managing-researchers-traders-and-engineers:1}
The new futures signal costs \$1.5 million, succeeds with probability 0.30 and is then worth \$12 million. What is its expected net value, and per
dollar of cost?
\end{exercise}

\begin{exercise}[$\star$]\label{exo:fm:managing-researchers-traders-and-engineers:2}
The second momentum variant costs \$1.0 million, succeeds with probability 0.40, is worth \$6 million and overlaps the book by 70\%. Should it be
funded?
\end{exercise}

\begin{exercise}[$\star$]\label{exo:fm:managing-researchers-traders-and-engineers:3}
Check that the Shapley credits of the table add up to the combined P\&L, and compute what leave-one-out leaves unallocated.
\end{exercise}

\begin{exercise}[$\star\star$]\label{exo:fm:managing-researchers-traders-and-engineers:4}
Two contributors bring identical signals worth 20 alone and 20 together; a third brings an independent signal worth 10. Compute the leave-one-out
and Shapley credits.
\end{exercise}

\begin{exercise}[$\star\star$]\label{exo:fm:managing-researchers-traders-and-engineers:5}
From \cref{fig:fm:managing-researchers-traders-and-engineers:dsr}, after how many variants does the chapter's backtest fall below an even chance of
beating luck?
\end{exercise}

\begin{exercise}[$\star\star$]\label{exo:fm:managing-researchers-traders-and-engineers:6}
Which credit rule rewards a researcher for replicating a signal the desk already trades, and why is that valuable?
\end{exercise}

\begin{exercise}[$\star\star\star$]\label{exo:fm:managing-researchers-traders-and-engineers:7}
\emph{Coding.} Rerun \texttt{fm\_research.credit(seed=12)}. How do the leave-one-out and Shapley credits change, and what does that say about
paying on one year's attribution?
\end{exercise}

\begin{exercise}[$\star\star\star$]\label{exo:fm:managing-researchers-traders-and-engineers:8}
\emph{Find the flaw.} ``We pay each researcher 10\% of what the desk would lose if their signal were switched off. That is fair: it is exactly
what they add.''
\end{exercise}

\section{Problem: Who Made the Money}

\begin{problem}[{Weekend problem --- who made the money}]\label{pb:fm:managing-researchers-traders-and-engineers:1}
A combined forecast made \$30 million. Three researchers ask for their share of a \$6 million pool, and the head of desk must also fund next year's
research.

\medskip\noindent\textbf{Part I --- The portfolio.}
\begin{enumerate}
\item Define a \omterm{def:fm:managing-researchers-traders-and-engineers:portfolio}{research portfolio} and a \omterm{def:fm:managing-researchers-traders-and-engineers:portfolio}{project scorecard}.
\item Which six projects does the greedy rule fund, and at what expected net value?
\item Give the mean realised value and the probability of a losing year for the greedy and the random choices.
\item Where do the probabilities of success on a scorecard come from, and why are they optimistic?
\end{enumerate}

\medskip\noindent\textbf{Part II --- The review.}
\begin{enumerate}[resume]
\item What makes a review theatre, and what prevents it?
\item Give the probabilistic and the deflated Sharpe ratio of the five-year backtest after 40 variants.
\item What is the expected best Sharpe ratio of 40 worthless variants?
\item Why must a new hypothesis start a new trial count?
\end{enumerate}

\medskip\noindent\textbf{Part III --- The credit.}
\begin{enumerate}[resume]
\item Define \omterm{def:fm:managing-researchers-traders-and-engineers:credit}{leave-one-out contribution} and \omterm{def:fm:managing-researchers-traders-and-engineers:credit}{credit attribution}.
\item Give each researcher's credit alone, by leave-one-out, by order of arrival and by Shapley.
\item State \cref{prop:fm:managing-researchers-traders-and-engineers:shapley} and prove its first part.
\item Why does leave-one-out leave most of the profit unallocated here?
\item Split the \$6 million pool by leave-one-out, Shapley and order of arrival.
\item How uncertain are these credits?
\end{enumerate}

\medskip\noindent\textbf{Part IV --- The people.}
\begin{enumerate}[resume]
\item What does each credit rule encourage researchers to do?
\item How should traders and engineers be judged while their results are noise?
\item Why do engineers belong in a strategy's \omterm{def:fm:managing-researchers-traders-and-engineers:credit}{credit attribution}?
\item Who should review a signal, and who should not?
\item State the \emph{named result}: the Shapley and leave-one-out credits of the three researchers, and the pool split each implies.
\item In two sentences, explain to the three researchers how the pool was split.
\end{enumerate}
\end{problem}

\section{Interview questions}

\begin{interviewq}[$\star$ \iqroles{researcher}]\label{iq:fm:managing-researchers-traders-and-engineers:1}
Two of your signals are highly correlated. What happens to each one's \omterm{def:fm:managing-researchers-traders-and-engineers:credit}{leave-one-out contribution} to the combined forecast?
\end{interviewq}

\begin{interviewq}[$\star$ \iqroles{researcher}]\label{iq:fm:managing-researchers-traders-and-engineers:2}
Your backtest has a Sharpe ratio of 1.2 over five years. What else do you need to say before anyone should believe it?
\end{interviewq}

\begin{interviewq}[$\star\star$ \iqroles{researcher, trader}]\label{iq:fm:managing-researchers-traders-and-engineers:3}
Define the Shapley value and give one property that leave-one-out credit lacks.
\end{interviewq}

\begin{interviewq}[$\star\star$ \iqroles{developer}]\label{iq:fm:managing-researchers-traders-and-engineers:4}
How would you measure an engineering team's contribution to a trading desk?
\end{interviewq}

\begin{interviewq}[$\star\star$ \iqroles{researcher}]\label{iq:fm:managing-researchers-traders-and-engineers:5}
You have a budget for five projects out of twelve. How do you choose, and which input do you trust least?
\end{interviewq}

\begin{interviewq}[$\star\star\star$ \iqroles{researcher}]\label{iq:fm:managing-researchers-traders-and-engineers:6}
Computing exact Shapley values needs $n!$ orders. How would you attribute credit among 30 contributors?
\end{interviewq}
